\documentclass[UTF8,11pt]{ctexart}
\usepackage[a4paper,margin=24mm]{geometry}
\usepackage{amsmath,amssymb,amsthm,mathtools,mathrsfs}
\usepackage[all]{xy}
\usepackage{xurl,hyperref,enumitem}
\usepackage{tikz}
\usetikzlibrary{arrows.meta,cd,decorations.markings,calc,patterns}
\hypersetup{hidelinks,breaklinks=true}
\setlength{\emergencystretch}{3em}
\allowdisplaybreaks
\newtheorem*{theorem}{Theorem}
\newtheorem*{thm}{Theorem}
\newtheorem*{lemma}{Lemma}
\newtheorem*{proposition}{Proposition}
\newtheorem*{prop}{Proposition}
\newtheorem*{corollary}{Corollary}
\newtheorem*{cor}{Corollary}
\newtheorem*{claim}{Claim}
\theoremstyle{definition}
\newtheorem*{definition}{Definition}
\newtheorem*{defn}{Definition}
\newtheorem*{remark}{Remark}
\newtheorem*{example}{Example}
\newcommand{\R}{\mathbb R}
\newcommand{\C}{\mathbb C}
\newcommand{\Z}{\mathbb Z}
\newcommand{\Q}{\mathbb Q}
\newcommand{\N}{\mathbb N}
\newcommand{\F}{\mathbb F}
\newcommand{\D}{\mathbb D}
\newcommand{\bB}{\mathbb B}
\newcommand{\bC}{\mathbb C}
\newcommand{\bR}{\mathbb R}
\newcommand{\bZ}{\mathbb Z}
\newcommand{\bN}{\mathbb N}
\newcommand{\bQ}{\mathbb Q}
\newcommand{\bD}{\mathbb D}
\newcommand{\bF}{\mathbb F}
\newcommand{\CP}{\mathbb{CP}}
\newcommand{\RP}{\mathbb{RP}}
\newcommand{\sO}{\mathscr O}
\newcommand{\sI}{\mathscr I}
\newcommand{\sF}{\mathscr F}
\newcommand{\sG}{\mathscr G}
\newcommand{\sH}{\mathscr H}
\newcommand{\sR}{\mathscr R}
\newcommand{\sM}{\mathscr M}
\newcommand{\sV}{\mathscr V}
\newcommand{\sU}{\mathscr U}
\DeclareMathOperator{\Hom}{Hom}
\DeclareMathOperator{\Ext}{Ext}
\DeclareMathOperator{\Tor}{Tor}
\DeclareMathOperator{\Spec}{Spec}
\DeclareMathOperator{\Proj}{Proj}
\DeclareMathOperator{\supp}{supp}
\DeclareMathOperator{\im}{im}
\DeclareMathOperator{\id}{id}
\DeclareMathOperator{\Tr}{Tr}
\DeclareMathOperator{\tr}{tr}
\DeclareMathOperator{\rank}{rank}
\DeclareMathOperator{\ord}{ord}
\DeclareMathOperator{\dist}{dist}
\DeclareMathOperator{\End}{End}
\DeclareMathOperator{\Aut}{Aut}
\DeclareMathOperator{\Gal}{Gal}
\DeclareMathOperator{\vol}{vol}
\DeclareMathOperator{\Cl}{Cl}
\DeclareMathOperator{\coker}{coker}
\DeclareMathOperator{\codim}{codim}
\DeclareMathOperator{\divergence}{div}
\DeclareMathOperator{\Ric}{Ric}
\DeclareMathOperator{\grad}{grad}
\DeclareMathOperator{\Hess}{Hess}
\newcommand{\abs}[1]{\left\lvert#1\right\rvert}
\newcommand{\sabs}[1]{\lvert#1\rvert}
\newcommand{\babs}[1]{\bigl\lvert#1\bigr\rvert}
\newcommand{\Babs}[1]{\Bigl\lvert#1\Bigr\rvert}
\newcommand{\bbabs}[1]{\biggl\lvert#1\biggr\rvert}
\newcommand{\BBabs}[1]{\Biggl\lvert#1\Biggr\rvert}
\newcommand{\norm}[1]{\left\lVert#1\right\rVert}
\newcommand{\snorm}[1]{\lVert#1\rVert}
\newcommand{\bnorm}[1]{\bigl\lVert#1\bigr\rVert}
\newcommand{\Bnorm}[1]{\Bigl\lVert#1\Bigr\rVert}
\newcommand{\bbnorm}[1]{\biggl\lVert#1\biggr\rVert}
\newcommand{\BBnorm}[1]{\Biggl\lVert#1\Biggr\rVert}
\newcommand{\widebar}[1]{\overline{#1}}
\newcommand{\nicefrac}[2]{\frac{#1}{#2}}
\newcommand{\linnprod}[2]{\langle#1,#2\rangle}
\newcommand{\myindex}[1]{#1}
\newcommand{\glsadd}[1]{}
\newcommand{\avoidbreak}{}
\newcommand{\sourceref}[1]{\textnormal{[原文：\nolinkurl{#1}]}}
\newcommand{\thmref}[1]{\sourceref{#1}}
\newcommand{\lemmaref}[1]{\sourceref{#1}}
\newcommand{\propref}[1]{\sourceref{#1}}
\newcommand{\corref}[1]{\sourceref{#1}}
\newcommand{\defnref}[1]{\sourceref{#1}}
\newcommand{\exampleref}[1]{\sourceref{#1}}
\newcommand{\exerciseref}[1]{\sourceref{#1}}
\newcommand{\figureref}[1]{\sourceref{#1}}
\newcommand{\sectionref}[1]{\sourceref{#1}}
\newcommand{\appendixref}[1]{\sourceref{#1}}
\newcommand{\stacksref}[2]{\href{https://stacks.math.columbia.edu/tag/#1}{#2 [#1]}}

\begin{document}
\section*{045\quad Riemann完备zeta函数的延拓与函数方程}
\subsection*{来源与计数目标}
Andrew V. Sutherland，MIT 18.785 \emph{Number Theory I}，Fall 2021，Lecture 17。\S17.1：Theorem 17.8、Lemmas 17.9--17.10（p.~3），\S17.1.3 的完整推导至式(3)--(4)（pp.~5--6）。采用 MIT OpenCourseWare 实际访问版本；获取日期：2026-09-28。
\par\url{https://ocw.mit.edu/courses/18-785-number-theory-i-fall-2021/mit18_785f21_lec17.pdf}
\par\textbf{本文件只计一个问题：}构造 $Z(s)=\pi^{-s/2}\Gamma(s/2)\zeta(s)$ 的全平面亚纯延拓并证明 $Z(s)=Z(1-s)$ 及两个简单极点；不增加零点分布要求。
\subsection*{作者命题与人类证明}
\paragraph{Notation (Definitions 17.1, 17.4, 17.11--17.12).}
A Schwartz function on $\R$ is a complex-valued $C^\infty$ function $f$ such that $\sup_{x\in\R}|x^m f^{(n)}(x)|<\infty$ for all $m,n\in\Z_{\geq0}$. Write $\mathcal S(\R)$ for this space and
\[
\widehat f(y)=\int_\R f(x)e^{-2\pi ixy}\,dx,\qquad
\Gamma(s)=\int_0^\infty e^{-t}t^{s-1}\,dt\quad(\operatorname{Re}s>0).
\]
\begin{theorem}[17.8 (Poisson Summation Formula)]
For all $f\in\mathcal S(\R)$ we have the identity
\[
\sum_{n\in\Z}f(n)=\sum_{n\in\Z}\widehat f(n).
\]
\end{theorem}
\begin{proof}
We first note that both sums are well defined; the rapid decay property of Schwartz functions guarantees absolute convergence. Let $F(x):=\sum_{n\in\Z}f(x+n)$. Then $F$ is a periodic $C^\infty$-function, so it has a Fourier series expansion
\[
F(x)=\sum_{n\in\Z}c_ne^{2\pi inx},
\]
with Fourier coefficients
\[
c_n=\int_0^1 F(t)e^{-2\pi int}\,dt
=\int_0^1\sum_{m\in\Z}f(t+m)e^{-2\pi int}\,dt
=\int_\R f(t)e^{-2\pi int}\,dt=\widehat f(n).
\]
We then note that $\sum_{n\in\Z}f(n)=F(0)=\sum_{n\in\Z}c_n=\sum_{n\in\Z}\widehat f(n)$.
\end{proof}
\begin{lemma}[17.9]
Let $g(x):=e^{-\pi x^2}$. Then $\widehat g(y)=g(y)$.
\end{lemma}
\begin{proof}
The function $g(x)$ satisfies the first order ordinary differential equation
\[
g'+2\pi xg=0,\tag{1}
\]
with initial value $g(0)=1$. Multiplying both sides by $-i$ and taking Fourier transforms yields
\[
-i(\widehat {g'}+2\pi\widehat{xg})=-i(2\pi ix\widehat g+i\widehat g')=\widehat g'+2\pi x\widehat g=0,
\]
via Lemma 17.7. So $\widehat g$ also satisfies (1), and $\widehat g(0)=\int_\R e^{-\pi x^2}\,dx=1$, so $\widehat g=g$.
\end{proof}
\paragraph{17.1.1 Jacobi's theta function}
We now define the theta function
\[
\Theta(\tau):=\sum_{n\in\Z}e^{\pi in^2\tau}.
\]
The sum is absolutely convergent for $\operatorname{im}\tau>0$ and thus defines a holomorphic function on the upper half plane. It is easy to see that $\Theta(\tau)$ is periodic modulo 2, that is, $\Theta(\tau+2)=\Theta(\tau)$, but it it also satisfies another functional equation.
\begin{lemma}[17.10]
For all $a\in\R_{>0}$ we have $\Theta(ia)=\Theta(i/a)/\sqrt a$.
\end{lemma}
\begin{proof}
Put $g(x):=e^{-\pi x^2}$ and $h(x):=g(\sqrt a x)=e^{-\pi x^2a}$. Lemmas 17.6 and 17.9 imply
\[
\widehat h(y)=\widehat{g(\sqrt a x)}(y)=\widehat g(y/\sqrt a)/\sqrt a=g(y/\sqrt a)/\sqrt a.
\]
Plugging $\tau=ia$ into $\Theta(\tau)$ and applying Poisson summation (Theorem 17.8) yields
\[
\Theta(ia)=\sum_{n\in\Z}e^{-\pi n^2a}=\sum_{n\in\Z}h(n)
=\sum_{n\in\Z}\widehat h(n)
=\sum_{n\in\Z}g(n/\sqrt a)/\sqrt a=\Theta(i/a)/\sqrt a.
\]
\end{proof}
\paragraph{17.1.3 Completing the zeta function}
\begin{proof}[Author's derivation preceding equations (3)--(4)]
Let us now consider the function $F(s):=\pi^{-s}\Gamma(s)\zeta(2s)$, which is holomorphic on $\operatorname{Re}(s)>1/2$. In the region $\operatorname{Re}(s)>1/2$ we have an absolutely convergent sum
\begin{align*}
F(s)&=\pi^{-s}\Gamma(s)\sum_{n\geq1}n^{-2s}
=\sum_{n\geq1}(\pi n^2)^{-s}\Gamma(s)\\
&=\sum_{n\geq1}\int_0^\infty(\pi n^2)^{-s}t^{s-1}e^{-t}\,dt,
\end{align*}
and the substitution $t=\pi n^2y$ with $dt=\pi n^2dy$ yields
\[
F(s)=\sum_{n\geq1}\int_0^\infty(\pi n^2)^{-s}(\pi n^2y)^{s-1}e^{-\pi n^2y}\pi n^2\,dy
=\sum_{n\geq1}\int_0^\infty y^{s-1}e^{-\pi n^2y}\,dy.
\]
By the Fubini--Tonelli theorem, we can swap the sum and the integral to obtain
\[
F(s)=\int_0^\infty y^{s-1}\sum_{n\geq1}e^{-\pi n^2y}\,dy.
\]
We have $\Theta(iy)=\sum_{n\in\Z}e^{-\pi n^2y}=1+2\sum_{n\geq1}e^{-\pi n^2y}$, thus
\begin{align*}
F(s)&=\frac12\int_0^\infty y^{s-1}(\Theta(iy)-1)\,dy\\
&=\frac12\left(\int_0^1y^{s-1}\Theta(iy)\,dy-\frac1s+\int_1^\infty y^{s-1}(\Theta(iy)-1)\,dy\right).
\end{align*}
We now focus on the first integral on the RHS. The change of variable $t=1/y$ yields
\[
\int_0^1y^{s-1}\Theta(iy)\,dy
=\int_\infty^1 t^{1-s}\Theta(i/t)(-t^{-2})\,dt
=\int_1^\infty t^{-s-1}\Theta(i/t)\,dt.
\]
By Lemma 17.10, $\Theta(i/t)=\sqrt t\Theta(it)$, and adding $-\int_1^\infty t^{-s-1/2}\,dt+\int_1^\infty t^{-s-1/2}\,dt=0$ yields
\begin{align*}
&=\int_1^\infty t^{-s-1/2}(\Theta(it)-1)\,dt+\int_1^\infty t^{-s-1/2}\,dt\\
&=\int_1^\infty t^{-s-1/2}(\Theta(it)-1)\,dt-\frac1{1/2-s}.
\end{align*}
Plugging this back into our equation for $F(s)$ we obtain the identity
\[
F(s)=\frac12\int_1^\infty(y^{s-1}+y^{-s-1/2})(\Theta(iy)-1)\,dy-\frac1{2s}-\frac1{1-2s},
\]
valid on $\operatorname{Re}(s)>1/2$. We now observe that $F(s)=F(\frac12-s)$ for $s\ne0,\frac12$, which allows us to analytically extend $F(s)$ to a meromorphic function on $\C$ with poles only at $s=0,1/2$. Replacing $s$ with $s/2$ leads us to define the completed zeta function
\[
Z(s):=\pi^{-s/2}\Gamma(\tfrac s2)\zeta(s),\tag{3}
\]
which is meromorphic on $\C$ and satisfies the functional equation
\[
Z(s)=Z(1-s).\tag{4}
\]
It has simple poles at 0 and 1 (and no other poles).
\end{proof}
\subsection*{整理说明（非作者正文）}
开头记号段据 Definitions 17.1、17.4、17.11--17.12 摘编，非作者连续逐字段落。其后的 Poisson 求和、Gaussian Fourier 不变性、theta 变换和 Mellin 分割推导按原文顺序收录；最后一个 proof 环境是对原文非编号推导的排版包裹，不是新证明。不把辅助结果另计题数。

标准先修为 Schwartz 函数 Fourier 变换的伸缩与微分性质（17.6--17.7）、光滑周期函数 Fourier 展开、Gaussian 积分、Gamma 的定义、Fubini--Tonelli 与参数全纯积分。原作者的 Gamma 反射公式及零点讨论不属于当前题面，未收录；没有把 theta 功能方程当作未经展开的核心黑箱。p.~6 图像已查看，p.~3 和 p.~5 的截图接口失败，但解析正文已取得；未声称逐页图像核对。难度依据是把算术 Dirichlet 级数经 Mellin 变换接到 Fourier 对称性，再分离极点并解析延拓，而非仅验证一个函数恒等式。
\subsection*{可修改的简短标注（非作者正文）}
$c$：zeta 函数的解析延拓与对称性

$a$：Fourier 变换；Poisson 求和；Gaussian 自对偶；theta 级数；Mellin 变换；积分分割；极点分离

\texttt{cross\_theory\_status = yes}。算术求和经 Gamma 积分变成 theta 的 Mellin 变换；Poisson 求和产生尺度倒置对称性，变换回去给出 zeta 的函数方程和延拓。位置：17.10 与 \S17.1.3，pp.~3、5--6。

全部标注均为非穷尽、可修改初稿；未标注不表示未使用。
\subsection*{许可与修改记录}
材料来自 MIT OpenCourseWare，作者与课程见来源。按该站所示 CC BY-NC-SA 4.0 许可复用，整理部分沿用同一许可。\url{https://creativecommons.org/licenses/by-nc-sa/4.0/}。2026-09-28：助手转录题证、调整数学排版并加入来源、范围和初稿标注；不暗示作者或 MIT 认可本次整理。所留为本次题证转录摘录，不是原 PDF 下载件。
\end{document}
